\documentclass[]{article}
\usepackage{lmodern}
\usepackage{amssymb,amsmath}
\usepackage{ifxetex,ifluatex}
\usepackage{fixltx2e} % provides \textsubscript
\ifnum 0\ifxetex 1\fi\ifluatex 1\fi=0 % if pdftex
  \usepackage[T1]{fontenc}
  \usepackage[utf8]{inputenc}
\else % if luatex or xelatex
  \ifxetex
    \usepackage{mathspec}
  \else
    \usepackage{fontspec}
  \fi
  \defaultfontfeatures{Ligatures=TeX,Scale=MatchLowercase}
\fi
% use upquote if available, for straight quotes in verbatim environments
\IfFileExists{upquote.sty}{\usepackage{upquote}}{}
% use microtype if available
\IfFileExists{microtype.sty}{%
\usepackage{microtype}
\UseMicrotypeSet[protrusion]{basicmath} % disable protrusion for tt fonts
}{}
\usepackage[margin=1in]{geometry}
\usepackage{hyperref}
\hypersetup{unicode=true,
            pdftitle={How Legislators Actually Invest Their Time},
            pdfborder={0 0 0},
            breaklinks=true}
\urlstyle{same}  % don't use monospace font for urls
\usepackage{longtable,booktabs}
\usepackage{graphicx,grffile}
\makeatletter
\def\maxwidth{\ifdim\Gin@nat@width>\linewidth\linewidth\else\Gin@nat@width\fi}
\def\maxheight{\ifdim\Gin@nat@height>\textheight\textheight\else\Gin@nat@height\fi}
\makeatother
% Scale images if necessary, so that they will not overflow the page
% margins by default, and it is still possible to overwrite the defaults
% using explicit options in \includegraphics[width, height, ...]{}
\setkeys{Gin}{width=\maxwidth,height=\maxheight,keepaspectratio}
\IfFileExists{parskip.sty}{%
\usepackage{parskip}
}{% else
\setlength{\parindent}{0pt}
\setlength{\parskip}{6pt plus 2pt minus 1pt}
}
\setlength{\emergencystretch}{3em}  % prevent overfull lines
\providecommand{\tightlist}{%
  \setlength{\itemsep}{0pt}\setlength{\parskip}{0pt}}
\setcounter{secnumdepth}{0}
% Redefines (sub)paragraphs to behave more like sections
\ifx\paragraph\undefined\else
\let\oldparagraph\paragraph
\renewcommand{\paragraph}[1]{\oldparagraph{#1}\mbox{}}
\fi
\ifx\subparagraph\undefined\else
\let\oldsubparagraph\subparagraph
\renewcommand{\subparagraph}[1]{\oldsubparagraph{#1}\mbox{}}
\fi

%%% Use protect on footnotes to avoid problems with footnotes in titles
\let\rmarkdownfootnote\footnote%
\def\footnote{\protect\rmarkdownfootnote}

%%% Change title format to be more compact
\usepackage{titling}

% Create subtitle command for use in maketitle
\newcommand{\subtitle}[1]{
  \posttitle{
    \begin{center}\large#1\end{center}
    }
}

\setlength{\droptitle}{-2em}

  \title{How Legislators Actually Invest Their Time}
    \pretitle{\vspace{\droptitle}\centering\huge}
  \posttitle{\par}
    \author{}
    \preauthor{}\postauthor{}
    \date{}
    \predate{}\postdate{}
  
\usepackage{graphicx}

\begin{document}
\maketitle

Outline for APSA paper

Representation in Congress is based, in part, on legislators' ability to
assist their constituents with the federal bureaucracy. Constituency
service is central to several definitions of representation or a key
input to explain the incumbency advantage and yet, too little is known
about when and how elected officials contact the bureaucracy on behalf
of their constituents. We use a massive new database of Congressional
contact with bureaucracies obtained through over 200 FOIA requests to
examine when legislators contact agencies, who contact those agencies,
and the purpose of the contact. Using this new data set, we show that
legislators who are more powerful in Washington advocate more for their
constituents, rather than less. This suggests that perceived choices
constituents may face between influential representatives and attentive
representatives is a false choice.

Introduction Since the start of the Republic, members of Congress have
provided constituency service to their constituents, helping channel and
articulate their requests to the government. Representation through
constituency service constitutes an important facet of legislators'
jobs. And its growth has been used to explain the presence of the
incumbency advantage, or used as a metric to assess how legislators
preferentially allocate their time to constituents.

Within many theories of representation, there has been a perceived
tradeoff between representation from powerful elected officials and
elected officials who are attentive to the district. As legislators
acquire power in the institution, it is often asserted that they obtain
``Potomac fever'' and dedicate less attention to constituents back in
the district. This similar argument is made in the popular press and XX.

While it holds intuitive appeal, we show that there is no trade-off
between representation from powerful legislators and attentive
legislators. Using a massive new collection of Congressional
correspondence with bureaucratic agencies obtained through over 200 FOIA
requests, we show that as legislators acquire more power they advocate
more for their constituents. This difference is apparent both
cross-sectionally---more powerful legislators issue more requests than
their less powerful colleagues---and within legislator---as legislators
acquire more power they issue more requests. This occurs, we argue,
because legislators use their additional power, in part, to solidify
their electoral base at home. To that end, we demonstrate ZZ.

To demonstrate how institutional power increases district attention, we
utilize a new data set of Congressional correspondence with the district
we obtained through over 200 FOIA requests. Using this new data set, we
also establish some important stylized facts about when and how elected
officials contact federal agencies. We show that MM, NN, OO.

We find that letter writing is not evenly distributed. Some members
simply use this tactic more than others, even within strata of
institutional power.

Furthermore, those who use this tactic use it broadly across issue areas
and parts of the government. The conventional model, where oversight
committee chairs are the most attentive to the parts of government under
jurisdiction largely holds in the House, where few rank and file members
regularly lobby the bureaucracy, but less so the Senate, where many
members are more active and often more attentive oversight committee
chairs.

Our results provide important new insights into how political
representation occurs in American politics. All things equal, we show
that constituents have a good reason to prefer more powerful
representatives. This clarifies why increased power is useful to elected
officials and helps establish how legislators are able to build their
capacity to both influence policy while ensuring that they maintain
support in their constituency.

\section{The False Power, District Attentiveness Tradeoff}

\begin{enumerate}
\def\labelenumi{\arabic{enumi})}
\tightlist
\item
  Who argues that there is this tradeoff?
\end{enumerate}

\begin{enumerate}
\def\labelenumi{\alph{enumi})}
\tightlist
\item
  Fenno,Mayhew
\item
  More recent articles (perhaps)
\end{enumerate}

\begin{enumerate}
\def\labelenumi{\arabic{enumi})}
\setcounter{enumi}{1}
\tightlist
\item
  Why might we expect this not to be true?
\end{enumerate}

\begin{enumerate}
\def\labelenumi{\alph{enumi})}
\tightlist
\item
  Level differences, based on ability
\item
  How individuals distribute power
\end{enumerate}

\section{Data: FOIA Requests}

Details on the FOIA requests 1) Agencies we've contacted: We requested
records from all cabinet departments, their subagencies, and 33
independent agencies. 2) Agencies we've heard from: As of August 2018,
all departments except for the State Department and Veterans' Affairs
have provided records, though the majority of Department of Defense
records are still being processed and the majority of Department of
Justice subagencies have either denied maintaining records or not yet
granted our request. As for independent agencies, we waiting on records
from the SEC, OPM, NRC, NSF, MSPB, FTC, and Export-Import Bank, CPSC,
CFTC, CIA, Appalachian Regional Commission. The remaining 22 independent
agencies have provided records, though some are still in the process of
reviewing and releasing additional records. Of these, 10 have been
cleaned and linked with other data source.

\begin{longtable}[]{@{}lccl@{}}
\toprule
\begin{minipage}[b]{0.18\columnwidth}\raggedright\strut
Department\strut
\end{minipage} & \begin{minipage}[b]{0.28\columnwidth}\centering\strut
Number of Agencies/Components\strut
\end{minipage} & \begin{minipage}[b]{0.27\columnwidth}\centering\strut
Number for which we have data\strut
\end{minipage} & \begin{minipage}[b]{0.16\columnwidth}\raggedright\strut
Number in Dataset\strut
\end{minipage}\tabularnewline
\midrule
\endhead
\begin{minipage}[t]{0.18\columnwidth}\raggedright\strut
Agriculture\strut
\end{minipage} & \begin{minipage}[t]{0.28\columnwidth}\centering\strut
25\strut
\end{minipage} & \begin{minipage}[t]{0.27\columnwidth}\centering\strut
18\strut
\end{minipage} & \begin{minipage}[t]{0.16\columnwidth}\raggedright\strut
18\strut
\end{minipage}\tabularnewline
\begin{minipage}[t]{0.18\columnwidth}\raggedright\strut
\ldots{}\strut
\end{minipage} & \begin{minipage}[t]{0.28\columnwidth}\centering\strut
\strut
\end{minipage} & \begin{minipage}[t]{0.27\columnwidth}\centering\strut
\strut
\end{minipage} & \begin{minipage}[t]{0.16\columnwidth}\raggedright\strut
\strut
\end{minipage}\tabularnewline
\begin{minipage}[t]{0.18\columnwidth}\raggedright\strut
Independent Agencies\strut
\end{minipage} & \begin{minipage}[t]{0.28\columnwidth}\centering\strut
33\strut
\end{minipage} & \begin{minipage}[t]{0.27\columnwidth}\centering\strut
22\strut
\end{minipage} & \begin{minipage}[t]{0.16\columnwidth}\raggedright\strut
10\strut
\end{minipage}\tabularnewline
\begin{minipage}[t]{0.18\columnwidth}\raggedright\strut
---------------------\strut
\end{minipage} & \begin{minipage}[t]{0.28\columnwidth}\centering\strut
:------------------------------:\strut
\end{minipage} & \begin{minipage}[t]{0.27\columnwidth}\centering\strut
:-----------------------------:\strut
\end{minipage} & \begin{minipage}[t]{0.16\columnwidth}\raggedright\strut
:-----------------\strut
\end{minipage}\tabularnewline
\begin{minipage}[t]{0.18\columnwidth}\raggedright\strut
TOTAL\strut
\end{minipage} & \begin{minipage}[t]{0.28\columnwidth}\centering\strut
\strut
\end{minipage} & \begin{minipage}[t]{0.27\columnwidth}\centering\strut
\strut
\end{minipage} & \begin{minipage}[t]{0.16\columnwidth}\raggedright\strut
\strut
\end{minipage}\tabularnewline
\bottomrule
\end{longtable}

What does a typical response look like?

What do we do with it then?

\subsection{Basic Facts About Contacts}

How many contacts do representatives make?

Who tends to make those contacts?\\
a) Early career? b) Ideological orientation? - No. See ideology section
in Summary c) Women more than men? Black vs Latinx vs other demographic
minorities?

What is the overall distribution look like (inequality) - Gini
coefficients are high. The Senate and House Gini coefficients are
comparable to income inequality in the US and Mexico - mean (per member
per year). With our current data, the mean correspondence per year is 18
for the house and 52 for the Senate. - means by type
\includegraphics{means_by_type_table} - density - density by type

Are there more contacts right before an election, right after or
somewhere in between?

In the House, both parties average just slightly more letters per year
while in the majority, but in the Senate, Republicans average more per
year while in the minority. \includegraphics{party_mean_by_month}

\subsection{More Power and More Contact}

We want to show this a few ways

\begin{enumerate}
\def\labelenumi{\alph{enumi})}
\tightlist
\item
  Show that there are level differences. Individuals with more power
  contact more. We can run this regression in several ways, including:
\item
  Committee chairs
\item
  Prestige Committees
\item
  Leadership positions within the party.
\end{enumerate}

Using those same variables (and other combinations) we can then run some
diff-in-diff specifications

The core model will be something like

Contacts in Year t \textasciitilde{} Legislator Fixed Effect + Congress
Fixed Effect + Legislator's Power in Year T + Time-varying covariates

Using this model we have a diff-in-diff where the relevant coefficient
is the legislator's power in year t. I``m happy to run these models.


\end{document}
